\section{Background and Related Work}

\subsection{The Efficiency Gap}

The efficiency gap measures partisan bias by calculating wasted votes for each party: votes cast for losing candidates and surplus votes beyond 50\%+1 needed to win~\cite{stephanopoulos2015partisan}. The metric ranges from $-50\%$ (complete Democratic advantage) to $+50\%$ (complete Republican advantage), with 0\% indicating no partisan bias. Efficiency gaps exceeding $\pm 7-8\%$ are generally considered indicative of extreme partisan gerrymandering~\cite{mcghee2020fairness}.

\subsection{Geographic Sorting and Unintentional Gerrymandering}

Rodden~\cite{rodden2019cities} and Chen \& Rodden~\cite{chen2013unintentional} demonstrate that Democratic voters' urban concentration creates structural disadvantages in single-member district systems. Even neutral redistricting produces efficiency gaps favoring Republicans because compact urban districts pack Democrats while sprawling rural districts efficiently allocate Republican votes.

\subsection{Algorithmic Redistricting}

Recent work has shown algorithmic methods can satisfy constitutional requirements~\cite{deluca2026recursive,cho2016measuring} and Voting Rights Act compliance~\cite{deluca2026vra}. However, partisan fairness analysis has been limited.

\subsection{Contribution}

This paper provides the first comprehensive efficiency gap analysis of algorithmic redistricting at national scale, establishing empirical benchmarks for partisan fairness evaluation.
